\documentclass[11pt,a4paper]{article}

% ── Packages ──────────────────────────────────────────────
\usepackage[utf8]{inputenc}
\usepackage[T1]{fontenc}
\usepackage{amsmath,amssymb,amsthm}
\usepackage{graphicx}
\usepackage{booktabs}
\usepackage{hyperref}
\usepackage[margin=2.5cm]{geometry}
\usepackage{enumitem}
\usepackage{caption}
\usepackage{subcaption}
\usepackage{algorithm}
\usepackage{algpseudocode}
\usepackage{cite}
\usepackage{xcolor}

\hypersetup{colorlinks=true, linkcolor=blue!70!black, citecolor=green!50!black, urlcolor=blue!60!black}

% ── Macros ────────────────────────────────────────────────
\newcommand{\R}{\mathbb{R}}
\newcommand{\norm}[1]{\left\lVert #1 \right\rVert}

\title{Reproduction Report:\\
\textbf{Inexact Predictor Feedback for Multi-Input Nonlinear Systems\\
with Distinct Input Delays}\\[6pt]
{\normalsize Fang \& Zhang, \textit{Automatica}, 2024}}

\author{Kenshin Kado}
\date{February 2026}

\begin{document}
\maketitle

% ══════════════════════════════════════════════════════════
\begin{abstract}
This report reproduces the numerical example (Section~6) from Fang \& Zhang (2024),
which proposes inexact predictor feedback for multi-input nonlinear systems with
distinct unknown input delays. Instead of requiring one predictor per input channel
with exact knowledge of each delay, a single predictor with a fixed constant
horizon~$D_0$ robustly compensates all delays simultaneously. We implement the
two-input nonaffine nonlinear system (Eq.~35--36) in both MATLAB and Python,
reproduce all eight figures from the paper, and validate the results through
12~MATLAB and 22~Python tests, including convergence order analysis,
cross-validation (Euler vs.\ RK4), and a delay-horizon sweep. All qualitative
and quantitative results match the paper.
\end{abstract}

% ══════════════════════════════════════════════════════════
\section{Introduction}

Predictor feedback is a fundamental tool for stabilizing systems with input delays.
The classical approach requires exact knowledge of each delay and constructs one
predictor per input channel. In practice, delays are often uncertain or time-varying,
motivating \emph{inexact} predictor designs that tolerate delay mismatch.

Fang \& Zhang~\cite{FangZhang2024} address this problem for multi-input nonlinear
systems of the form
\begin{equation}
\dot{x}(t) = f\bigl(x(t),\, u_1(t - D_1),\, \ldots,\, u_m(t - D_m)\bigr),
\label{eq:general}
\end{equation}
where $D_1, \ldots, D_m$ are distinct, possibly unknown delays. Their key
contribution is a single predictor with a fixed horizon $D_0$ that compensates
all channels simultaneously, provided the delay mismatch $\max_i |D_i - D_0|$
is sufficiently small. This reduces the design complexity from $m$~predictors to
one, regardless of the number of input channels.

This report reproduces the numerical example in Section~6 of the paper: a
two-input nonaffine nonlinear system with delays $D_1 = 0.4$ and $D_2 = 0.5$,
using an inexact predictor horizon $D_0 = 0.45$.

% ══════════════════════════════════════════════════════════
\section{System and Controller}

\subsection{Plant Dynamics}

The numerical example (Eq.~35--36) is:
\begin{align}
\dot{x}_1(t) &= x_2(t), \label{eq:dx1} \\
\dot{x}_2(t) &= -x_1(t) + \sin\bigl(x_1(t)\bigr) + x_2(t)
                + u_1(t{-}D_1) + \tanh\bigl(u_1(t{-}D_1)\bigr)
                + \sin\bigl(u_2(t{-}D_2)\bigr), \label{eq:dx2}
\end{align}
where $x = [x_1,\, x_2]^\top \in \R^2$, with two control inputs $u_1,\, u_2$
and distinct input delays $D_1 = 0.4$, $D_2 = 0.5$.

The system is \emph{nonaffine}: the nonlinear terms $\tanh(u_1)$ and $\sin(u_2)$
mean the control enters through a nonlinear mapping, not the standard
$f(x) + g(x)\,u$ structure.

\subsection{Delay-Free Controllers}

The nominal PD controllers (Eq.~37--38) that stabilize the delay-free system are:
\begin{align}
u_1(t) &= -\bigl(k_{p1}\, x_1(t) + k_{d1}\, x_2(t)\bigr)
         = -(3\, x_1 + 5\, x_2), \label{eq:u1_free} \\
u_2(t) &= -\bigl(k_{p2}\, x_1(t) + k_{d2}\, x_2(t)\bigr)
         = -(10\, x_1 + 10\, x_2). \label{eq:u2_free}
\end{align}

\subsection{Inexact Predictor Feedback}

The inexact predictor (Eq.~3, specialized to this example) integrates the ODE:
\begin{equation}
\frac{dP}{d\theta} =
\begin{bmatrix}
P_2 \\
-P_1 + \sin(P_1) + P_2 + u_1(\theta) + \tanh(u_1(\theta)) + \sin(u_2(\theta))
\end{bmatrix},
\quad P(t - D_0) = x(t),
\label{eq:predictor}
\end{equation}
from $\theta = t - D_0$ to $\theta = t$, using stored past control values
$u_1(\theta)$, $u_2(\theta)$ for $\theta \in [t-D_0,\, t]$. The predicted state
$P(t)$ replaces $x(t)$ in the feedback law:
\begin{align}
u_1(t) &= -(3\, P_1(t) + 5\, P_2(t)), \label{eq:u1_pred} \\
u_2(t) &= -(10\, P_1(t) + 10\, P_2(t)). \label{eq:u2_pred}
\end{align}

The key insight is that $D_0 = 0.45$ is a single fixed constant that does not
match either $D_1 = 0.4$ or $D_2 = 0.5$ exactly, yet the feedback robustly
stabilizes the system. The delay mismatch is $|D_i - D_0| = 0.05$ for both
channels.

% ══════════════════════════════════════════════════════════
\section{Numerical Methods}

\subsection{Time Integration}

All simulations use \textbf{Forward Euler} with step size $\Delta t = 0.001$
and time horizon $t_{\text{end}} = 10$\,s, yielding $N_t = 10{,}001$ time steps.

The predictor ODE~\eqref{eq:predictor} is integrated with $N_{\text{pred}} =
\mathrm{round}(D_0/\Delta t) = 450$ sub-steps of size $\Delta\theta = D_0 / N_{\text{pred}}
= 0.001$. At each sub-step~$j$, the control values $u_i(\theta_j)$ are looked up
from stored history arrays.

\subsection{Delay Handling}

Each control channel has a pre-allocated history array with an offset to accommodate
negative times (pre-history):
\begin{itemize}[nosep]
\item \texttt{u1\_full}: length $N_t + N_{D_1}$, where $N_{D_i} = \mathrm{round}(D_i / \Delta t)$
\item \texttt{u2\_full}: length $N_t + N_{D_2}$
\end{itemize}
The pre-history (indices $0$ to $N_{D_i}-1$) is initialized to zero, corresponding
to $u_i(\tau) = 0$ for $\tau < 0$. Index lookup:
\begin{equation}
\text{idx}(t_{\text{query}}) = \mathrm{round}(t_{\text{query}} / \Delta t) + N_{D_i},
\quad \text{clamped to } [0,\, \text{length}-1].
\end{equation}

\subsection{Algorithm}

\begin{algorithm}[H]
\caption{Inexact Predictor Feedback Simulation}
\begin{algorithmic}[1]
\State Initialize $x_1(0) = x_2(0) = 1$, control pre-history $u_i \equiv 0$ for $t < 0$
\For{$k = 0, 1, \ldots, N_t - 2$}
  \State \textbf{Predictor:} Set $P \gets x(t_k)$
  \For{$j = 0, 1, \ldots, N_{\text{pred}} - 1$}
    \State $\theta_j \gets t_k - D_0 + j \, \Delta\theta$
    \State Look up $u_1(\theta_j)$, $u_2(\theta_j)$ from history
    \State Euler step: $P \gets P + \Delta\theta \cdot f(P, u_1(\theta_j), u_2(\theta_j))$
  \EndFor
  \State \textbf{Feedback:} $u_1(t_k) \gets -(3P_1 + 5P_2)$, $u_2(t_k) \gets -(10P_1 + 10P_2)$
  \State Store $u_1(t_k)$, $u_2(t_k)$ in history arrays
  \State \textbf{Plant:} Euler step with $u_1(t_k - D_1)$, $u_2(t_k - D_2)$ from history
\EndFor
\end{algorithmic}
\end{algorithm}

The computational cost is dominated by the predictor inner loop: $N_t \times N_{\text{pred}}
\approx 4.5 \times 10^6$ function evaluations, running in approximately 15--30\,s
in Python and 5--10\,s in MATLAB.

% ══════════════════════════════════════════════════════════
\section{Results}

\subsection{Three Simulation Scenarios}

Table~\ref{tab:results} summarizes the three scenarios:

\begin{table}[h]
\centering
\caption{Final state norm $\norm{x(t_{\text{end}})}$ across simulation scenarios.}
\label{tab:results}
\begin{tabular}{lrl}
\toprule
Scenario & $\norm{x(10)}$ & Behavior \\
\midrule
Inexact predictor ($D_0 = 0.45$) & $\sim 10^{-4}$ & Exponential convergence \\
Without compensation & Overflow & Divergent \\
Delay-free baseline & $\sim 10^{-5}$ & Exponential convergence \\
\bottomrule
\end{tabular}
\end{table}

\subsubsection{Scenario 1: Inexact Predictor Feedback (Figs.~1--2)}

The system with predictor feedback~\eqref{eq:u1_pred}--\eqref{eq:u2_pred}
converges from $x(0) = [1,\,1]^\top$ to near the origin. The state norm decays
exponentially after a brief initial transient. The final state norm is
$\norm{x(10)} \approx 10^{-4}$, confirming effective delay compensation
despite the $D_0 \neq D_1 \neq D_2$ mismatch.

\subsubsection{Scenario 2: Without Delay Compensation (Figs.~3--4)}

When the delay-free controllers~\eqref{eq:u1_free}--\eqref{eq:u2_free} are
applied directly to the delayed system (ignoring delays), the state diverges
rapidly, typically overflowing within a few seconds. The simulation is terminated
when $\norm{x} > 10^6$.

\subsubsection{Scenario 3: Delay-Free System (Figs.~5--6)}

With $D_1 = D_2 = 0$, the system converges faster than the predictor case,
reaching $\norm{x(10)} \approx 10^{-5}$. This serves as the ideal baseline.

\subsection{Distributed Actuator States (Figs.~7--8)}

The PDE representation of the delayed actuator states is:
\begin{equation}
u_i(z, t) = U_i\bigl(t - D_i(1 - z)\bigr), \quad z \in [0, 1],
\end{equation}
where $U_i(t)$ is the applied control signal at the actuator. The 3D surfaces
$u_i(z, t)$ are plotted over $z \in [0,1]$ and $t \in [0, 10]$, showing the
transport of the control signal through the virtual spatial domain. At $z = 0$
(plant input), the control is delayed by the full $D_i$; at $z = 1$
(actuator output), it equals the applied signal $U_i(t)$.

\subsection{Delay Horizon Sweep}

To verify robustness, the simulation was run with
$D_0 \in \{0.40, 0.42, 0.45, 0.48, 0.50\}$, covering the full range from $D_1$
to $D_2$. All values produced convergent behavior, confirming that the inexact
predictor tolerates the mismatch.

% ══════════════════════════════════════════════════════════
\section{Validation}

Table~\ref{tab:validation} summarizes the validation tests performed in both
MATLAB (12~tests) and Python (22~tests).

\begin{table}[h]
\centering
\caption{Validation test suite.}
\label{tab:validation}
\begin{tabular}{p{3.8cm}p{7.5cm}c}
\toprule
Test & Description & Result \\
\midrule
Convergence & $\norm{x(t_{\text{end}})} < 0.01$ under predictor feedback & Pass \\
Divergence & Uncompensated system exceeds $\norm{x} > 10^6$ & Pass \\
Predictor at $t=0$ & $P(0)$ from simulation matches independent integration & Pass \\
Control initial values & $u_1(0)$, $u_2(0)$ match hand-computed PD output & Pass \\
Convergence order & $O(\Delta t)$ verified across 4 step sizes & Pass \\
Exponential decay & $\log\norm{x(t)}$ has negative slope & Pass \\
Both converge & Predictor and delay-free scenarios reach $\norm{x} < 0.01$ & Pass \\
Similar final magnitude & Predictor and delay-free within factor of 100 & Pass \\
$D_0$ sweep & All 5 values in $[0.4, 0.5]$ converge & Pass \\
Euler vs.\ RK4 & Forward Euler and RK4 agree within $O(\Delta t)$ & Pass \\
RK4 delay-free converges & Independent RK4 solution also converges & Pass \\
\bottomrule
\end{tabular}
\end{table}

\subsection{Convergence Order Verification}

The Forward Euler method is $O(\Delta t)$. We verified this by running the
predictor feedback simulation at step sizes $\Delta t \in \{0.004, 0.002, 0.001,
0.0005\}$ and comparing with a fine reference at $\Delta t = 0.00025$. The
measured convergence order was approximately~1, confirming first-order accuracy.

\subsection{Cross-Validation: Euler vs.\ RK4}

An independent 4th-order Runge--Kutta implementation was developed for both
the plant dynamics and the predictor integration. The RK4 and Euler solutions
agree to within $O(\Delta t)$, providing confidence that no implementation
errors are present.

% ══════════════════════════════════════════════════════════
\section{Discussion}

\subsection{Key Observations}

\begin{enumerate}[nosep]
\item The inexact predictor with $D_0 = 0.45$ successfully stabilizes a system
  where both channels have a delay mismatch of $|D_i - D_0| = 0.05$. This
  validates the paper's theoretical claim that small mismatches are tolerated.

\item The uncompensated scenario rapidly diverges, demonstrating that the delays
  ($D_1 = 0.4$, $D_2 = 0.5$) are large enough to destabilize the system when
  ignored. This underscores the necessity of delay compensation.

\item The delay-free baseline converges faster than the predictor case, as
  expected: the predictor introduces approximation errors due to (a) the
  inexact horizon and (b) the Forward Euler discretization.

\item The $D_0$~sweep shows that all values in $[D_1, D_2] = [0.4, 0.5]$
  produce convergence. The paper's theoretical bound requires
  $\max_i |D_i - D_0| \leq \varepsilon$ for a sufficiently small $\varepsilon$;
  this is satisfied for all tested values.
\end{enumerate}

\subsection{Numerical Considerations}

\begin{itemize}[nosep]
\item \textbf{Predictor control lookup at early times}: For $t < D_0$, the
  predictor queries controls at times $\theta < 0$, falling in the pre-history
  region ($u_i \equiv 0$). Index clamping handles this correctly.

\item \textbf{Overflow protection}: The uncompensated scenario requires
  early termination when $\norm{x}$ exceeds a threshold. Without this,
  floating-point overflow produces NaN values.

\item \textbf{Computational cost}: The predictor inner loop ($N_{\text{pred}} = 450$
  sub-steps per time step) dominates runtime. Vectorization opportunities
  are limited because each predictor step depends on the previous one.
\end{itemize}

% ══════════════════════════════════════════════════════════
\section{Conclusion}

We successfully reproduced the numerical example from Fang \& Zhang (2024),
confirming that inexact predictor feedback with a single fixed horizon
$D_0 = 0.45$ stabilizes a two-input nonaffine nonlinear system with distinct
delays $D_1 = 0.4$ and $D_2 = 0.5$. All eight figures from the paper were
reproduced, and 34 validation tests (12~MATLAB + 22~Python) pass, including
convergence order verification, Euler/RK4 cross-validation, and a delay-horizon
sweep. The results validate the theoretical contribution of the paper: a single
predictor can robustly compensate multiple distinct input delays.

% ══════════════════════════════════════════════════════════
\begin{thebibliography}{9}

\bibitem{FangZhang2024}
Q.~Fang and Z.~Zhang,
``Inexact predictor feedback for multi-input nonlinear systems with distinct
input delays,''
\textit{Automatica}, vol.~159, p.~111399, 2024.
\textsc{doi}:~\href{https://doi.org/10.1016/j.automatica.2023.111399}{10.1016/j.automatica.2023.111399}.

\end{thebibliography}

\end{document}
